\section{One conserved stress tensor on the comparison curve}
\label{sec:actual-stress-normalization}

The reference stress tensor must keep its prescribed sign as the requested
grades, derivative orders and configuration windows increase. We construct
one physical tensor and one comparison curve with that property. The only
subsequence taken below is a subsequence of the comparison parameter. The
physical correlation table, and the cutoff extraction that produced it,
remain fixed.

Write $E_d$ for the reference grade-$d$ field space and $F_d$ for the
physical one. Both have their given units, conjugations, signed coordinate
actions and inclusions between grades. Their derivative maps are
$\partial^0_{d,a}:E_d\to E_{d+1}$ and
$\partial^1_{d,a}:F_d\to F_{d+1}$, where $a\in\{0,1,2,3\}$.
Both correlation hierarchies are smooth, multilinear and faithful: a field
whose separated tested correlations with every spectator word vanish is
zero. Their units have expectation one and act as scalar insertions.
We use the exact derivative identities for separated correlations, with
derivative maps that annihilate units and respect reality and the signed
coordinate action. All these spaces are finite-dimensional. A comparison curve consists of
complex-linear isomorphisms $U_{k,d}:E_d\to F_d$ and positive radii
$r_k$ strictly decreasing to zero. For each fixed grade, the maps
eventually preserve the units, conjugations and signed actions; for
each fixed pair of grades they eventually commute with the inclusions.
These are the corrected comparison maps constructed in
Section~\ref{sec:ex-lattice-trace-producer}.

We retain the strength of their test estimate explicitly. If
$\boldsymbol v=(v_1,\ldots,v_n)$ is any fixed tuple of reference fields,
put
\[
 K_{k,\boldsymbol v}(x)
 =S^1_{U_k\boldsymbol v}(r_kx)-S^0_{\boldsymbol v}(x).
\]
For every compact $Q\subset\operatorname{Conf}_n(\mathbb R^4)$ and
$M>0$, the estimate is
\begin{equation}\label{eq:stress-uniform-tests}
 \sup_{\substack{f\in C_c^\infty,\ \operatorname{supp}f\subset Q\\
                         \|f\|_\infty\le M}}
       \left|\int f(x)K_{k,\boldsymbol v}(x)\,dx\right|
                 \longrightarrow0.
\end{equation}
The kernels are smooth on configuration space. The short-test bound
\eqref{eq:ex-lattice-trace-test-error} has precisely this uniformity:
its right side depends on $M$ and the labels, not on derivatives of $f$.
The finite slot expansions in the correction of the comparison maps
preserve it. Indeed, each expansion has finitely many fixed label tuples;
coefficients converging to zero multiply a fixed smooth kernel, whose
integral is bounded by $M|Q|$ times its supremum on $Q$; bounded
coefficients multiply errors satisfying \eqref{eq:stress-uniform-tests}.
Thus we use the uniform short-test receipt, not merely convergence for
each individual test.

\subsection{Derivative transport and convergence on compact sets}
Define the transported derivative
\begin{equation}\label{eq:stress-transported-derivative}
 A_{k,d,a}=r_k U_{k,d+1}^{-1}\partial^1_{d,a}U_{k,d}.
\end{equation}
It converges in operator norm to $\partial^0_{d,a}$. Here is the
finite-dimensional argument, including why its argument need not be
bounded in advance. Choose finitely many tested reference correlations
which detect $E_{d+1}$, as in the separating-test argument of
Section~\ref{sec:ex-lattice-trace-producer}. Let
$L:E_{d+1}\to\mathbb C^q$ be their injective measurement map and
$L_k$ the same measurements in the scaled physical table. On a finite
basis, \eqref{eq:stress-uniform-tests} gives $L_k\to L$ in operator
norm. Compactness of the unit sphere gives $c>0$ with
$\|Lv\|\ge c\|v\|$, and hence eventually
$\|L_kv\|\ge(c/2)\|v\|$.
For each fixed $v\in E_d$, the chain rule and the derivative identity
express each component of $L_k A_{k,d,a}v$ as a correlation tested
against minus the corresponding derivative of the original test.
Integration by parts has no boundary term because the test has compact
support inside configuration space. The approach estimate therefore gives
\[
 L_k A_{k,d,a}v\longrightarrow L\partial^0_{d,a}v.
\]
Subtract $L_k\partial^0_{d,a}v$, which has the same limit, and use the
lower bound for $L_k$. It follows that
$A_{k,d,a}v\to\partial^0_{d,a}v$. Checking a finite basis proves the
claimed operator-norm convergence.

The approach estimate also holds when each label $v_i(k)$ converges to
$v_i$. Expand the labels in fixed finite bases. The difference is a
finite sum of bounded coefficients multiplying fixed-tuple errors, plus
coefficients tending to zero multiplying fixed reference kernels.
The same compact integral bound proves
\eqref{eq:stress-uniform-tests} for this difference. Differentiating
one position replaces its label by $A_{k,d,a}v_i(k)$, which converges
to $\partial^0_{d,a}v_i$. Consequently this class of differences is
closed under any fixed number of position derivatives.

We give the compact convergence upgrade used here. Let $h_k$ be any
such difference, let $\chi$ be a smooth compactly supported cutoff
inside configuration space, and use the Fourier convention
$\widehat g(\xi)=\int e^{-2\pi i\xi\cdot x}g(x)\,dx$.
The tests $\chi(x)e^{-2\pi i\xi\cdot x}$ have one common support and
supremum bound. Equation~\eqref{eq:stress-uniform-tests} therefore yields
\[
       \sup_\xi|\widehat{\chi h_k}(\xi)|\longrightarrow0.
\]
For any fixed coordinate $a$ and integer $b\ge0$, integration by
parts and the product rule give
\[
 (2\pi i\xi_a)^b\widehat{\chi h_k}(\xi)
 =\sum_{j=0}^b {b\choose j}
       \widehat{(\partial_a^{b-j}\chi)(\partial_a^jh_k)}(\xi).
\]
Each term on the right tends uniformly to zero by derivative closure.
In the dimension $D=4n>0$, choose $b>D$ and a coordinate with
$|\xi_a|\ge|\xi|/\sqrt D$. Combining the finitely many coordinate
bounds with the unweighted bound gives numbers $\epsilon_k\to0$ such
that
\[
 |\widehat{\chi h_k}(\xi)|
       \le C_b\epsilon_k(1+|\xi|)^{-b}.
\]
The right side is integrable. Fourier inversion for the smooth compactly
supported function $\chi h_k$ gives
$\|\chi h_k\|_\infty\to0$. Given a compact set in configuration
space, choose $\chi=1$ on a neighborhood of it. Repeat the argument
with every coordinate derivative of a fixed order; finite dimensionality
then gives convergence in the corresponding derivative norm, uniformly
on the compact set. At arity zero the space of configurations is a
point and the assertion is the scalar approach estimate. We have proved
compact convergence of all fixed jets, including at coordinate ties that
are not collisions.

\subsection{The physical conserved line}
Let $K^0$ and $K^1$ be the real spaces of centered, symmetric,
real, signed-covariant grade-four stress arrays in the reference and
physical tables. Centering means that every entry has expectation zero.
Let $V^0,V^1$ be the spaces of real, signed-covariant grade-five vector
arrays, without imposing a centering condition, and
write
\[
 (D^iB)_a=\sum_b\partial^i_{4,b}B_{ba}\qquad(i=0,1)
\]
for divergence. Assume the fixed reference tensor $J^0$ belongs to
$K^0$, is nonzero and conserved, and that $D^0$ is onto with
$\ker D^0=\mathbb R J^0$. These are the reference rank and
normalization requirements; the low-grade rank calculation is
Section~\ref{sec:ex-lattice-trace-free-rank}. The actual reference
Ward normalization must be supplied for this same $J^0$.

For all sufficiently large $k$, the centered map
\[
 C_kB=\operatorname{center}_1(U_{k,4}B):K^0\longrightarrow K^1
\]
is an isomorphism. Its inverse is
$B\mapsto\operatorname{center}_0(U_{k,4}^{-1}B)$: the two scalar
unit corrections cancel exactly, by the calculation in
Section~\ref{sec:ex-lattice-trace-producer}. Unit preservation is
essential; the two expectation functionals need not agree under $U_k$.
The componentwise map $V_k:V^0\to V^1$ induced by $U_{k,5}$ is
also an isomorphism. Define
\[
       B_k=V_k^{-1}r_kD^1 C_k:K^0\longrightarrow V^0.
\]
Derivatives annihilate scalar units, so
$(B_kB)_a=\sum_b A_{k,4,b}B_{ba}$. Thus $B_k\to D^0$ in
operator norm. Choose a linear right inverse $R:V^0\to K^0$ of
the surjection $D^0$. Then $B_kR\to I$. For large $k$,
$\|B_kR-I\|<1/2$, so $B_kR$ is injective and hence bijective
in its finite-dimensional space. Therefore $B_k$ is onto.
Conjugating by $C_k,V_k$ and dividing by $r_k>0$ shows that $D^1$
is onto. The two pairs of spaces have equal dimensions, so
rank--nullity gives
\begin{equation}\label{eq:stress-physical-line}
                  \dim_{\mathbb R}\ker D^1=1.
\end{equation}

\subsection{Fixing the normalization and sign once}
Choose any $0\ne B\in\ker D^1$. In the full reference stress-array
space, including scalar directions, put
\[
 b_k=U_{k,4}^{-1}(r_k^4B),\qquad
 \xi_k=\frac{\|b_k\|}{\|J^0\|},\qquad
                       \widetilde B_k=\xi_k^{-1}b_k.
\]
Invertibility and $r_k>0$ imply $\xi_k>0$; no uniform lower bound
on $\xi_k$ is needed. We have
$\|\widetilde B_k\|=\|J^0\|$. Compactness of this sphere supplies
a strictly increasing subsequence on which $\widetilde B_k\to B_*$,
with $\|B_*\|=\|J^0\|$. Restrict $U_k,r_k,\xi_k$ to that
subsequence. All approach and eventual coherence statements persist.

The limit $B_*$ has the reference symmetry, reality and signed
covariance, since these are closed linear conditions. Its centering
requires a separate check. The corresponding physical entry is
\[
 U_{k,4}\widetilde B_{k,ab}=\xi_k^{-1}r_k^4B_{ab},
\]
whose expectation is exactly zero for every $k$. Apply the moving-label
compact convergence proved above at arity one and position zero.
It identifies the limit of these expectations with
$\mathbb E_0(B_{*,ab})$, proving that $B_*\in K^0$.
Moreover the exact physical conservation law gives
\[
 \sum_b A_{k,4,b}\widetilde B_{k,ba}
       =\xi_k^{-1}r_k^5 U_{k,5}^{-1}
                        \left(\sum_b\partial^1_{4,b}B_{ba}\right)=0.
\]
Operator-norm convergence of $A_{k,4,b}$ and convergence of the
balanced arrays imply $D^0B_*=0$. Thus $B_*=cJ^0$ for some real
$c$. Equality of the nonzero norms gives $|c|=1$.
Choose $J=B$ if $c=1$ and $J=-B$ if $c=-1$. This one choice gives
\begin{equation}\label{eq:stress-fixed-normalization}
 U_{k,4}^{-1}(\xi_k^{-1}r_k^4J_{ab})
                      \longrightarrow J^0_{ab}\qquad(a,b=0,1,2,3).
\end{equation}
The physical tensor $J$ is nonzero, centered and conserved. The
reference tensor has not been rescaled or changed to accommodate
later requests.

\subsection{All finite requests use tails of this curve}
Fix an arity bound $N$, a grade bound $d$, a derivative order $m$ and
$R\ge1$. For $n\le N$, use the compact window
\[
 Q_{n,R}=\{x:|x|\le R,\ |x_i-x_j|_\infty\ge R^{-1}
                                      \text{ for }i\ne j\}.
\]
Choose a single tail on which the exact symmetry, unit and inclusion
identities hold through grade $\max(d,4)$. There are only finitely
many grades and pairs involved. On this tail the compact convergence
argument gives, for every fixed tuple of grades at most $d$ and every
$j\le m$, convergence of its $j$th derivative uniformly on $Q_{n,R}$.
The tolerance threshold may depend on the tuple; no uniform bound over
all labels is asserted. Equation~\eqref{eq:stress-fixed-normalization}
holds on the same tail. Increasing the finite request only requires
another tail of this single curve, with the same $J,J^0$ and positive
normalization sequence. This proves the coherent comparison needed
before the annular Ward identities are applied. The uniform test
approach, exact derivative laws and reference conserved-line certificate
remain its stated inputs; the normalization argument does not supply
their analytic construction.
